% =====================================================================
%  quy gán nguồn cho rag hai tầng — slide
%  biên dịch:  build/tools/tectonic.exe slides.tex
%  ghép:       python build/mk_slides.py
% =====================================================================
\documentclass[aspectratio=169,10pt]{beamer}

\usepackage{fontspec}
\setmainfont{texgyrepagella}[
  Extension      = .otf,
  UprightFont    = *-regular,
  BoldFont       = *-bold,
  ItalicFont     = *-italic,
  BoldItalicFont = *-bolditalic ]
% Scale=1.0: o 0.82 khoi dia chi con 5,7pt — khong doc noi khi chieu len man.
\setmonofont{lmmono10-regular}[Extension=.otf, Scale=1.0]

\usepackage{tcolorbox}
\usepackage{fancyvrb}
\usepackage{amsmath, amssymb}
\usepackage{booktabs}
\usepackage{array}

% ── bảng màu ────────────────────────────────────────────────────────
\definecolor{navy}{HTML}{0E3F92}
\definecolor{tint}{HTML}{EDF2FB}
\definecolor{ink}{HTML}{2B2B2B}
\definecolor{mute}{HTML}{6B7280}
\definecolor{warn}{HTML}{9A3412}
\definecolor{good}{HTML}{14532D}
\definecolor{badge}{HTML}{A8C4EE}

% ── theme tối giản ─────────────────────────────────────────────────
\usetheme{default}
\usecolortheme{default}
\setbeamertemplate{navigation symbols}{}
\setbeamercolor{normal text}{fg=ink}
\setbeamercolor{structure}{fg=navy}
\setbeamercolor{frametitle}{fg=white, bg=navy}
\setbeamerfont{frametitle}{size=\large, series=\bfseries}

% Phu de khung = NHAN BUOC (a1..a5, b1..b3).  Template mac dinh cua beamer
% bi ta thay bang cai duoi, nen phai tu goi \insertframesubtitle — thieu no
% thi moi \framesubtitle am tham bien mat.
\setbeamertemplate{frametitle}{%
  \nointerlineskip
  \begin{beamercolorbox}[wd=\paperwidth, ht=2.4ex, dp=1.2ex,
                         leftskip=0.9em, rightskip=0.9em]{frametitle}
    \usebeamerfont{frametitle}\insertframetitle
    \hfill
    {\normalfont\footnotesize\color{badge}\insertframesubtitle}
  \end{beamercolorbox}}

\setbeamertemplate{footline}{%
  \hbox{\begin{beamercolorbox}[wd=\paperwidth, ht=2.2ex, dp=1ex]{}
    \hspace{0.9em}\color{mute}\scriptsize\insertsection
    \hfill\color{mute}\scriptsize\insertframenumber\,/\,\inserttotalframenumber
    \hspace{0.9em}\end{beamercolorbox}}}

\setbeamertemplate{itemize item}{\color{navy}\small$\bullet$}
\setbeamertemplate{itemize subitem}{\color{mute}\scriptsize$\circ$}
\setlength{\leftmargini}{1.1em}
\setlength{\leftmarginii}{1.1em}

% ── hộp ─────────────────────────────────────────────────────────────
\tcbset{
  % before/after skip: mac dinh la \medskipamount (~5pt) moi dau — tren slide
  % xep ba hop chong nhau thi rieng khoang trang da an 20pt chieu cao.
  bx/.style={colback=tint, colframe=navy!22, boxrule=0.4pt, arc=2pt,
             left=4.5pt, right=4.5pt, top=2pt, bottom=2pt,
             before skip=3pt, after skip=3pt,
             fonttitle=\bfseries\small, coltitle=navy},
  key/.style={bx, colback=navy!6, colframe=navy!45},
  bad/.style={bx, colback=warn!6, colframe=warn!35, coltitle=warn},
}
\newtcolorbox{bx}[1][]{bx,#1}
\newtcolorbox{keyb}[1][]{key,#1}
\newtcolorbox{badb}[1][]{bad,#1}

\DefineVerbatimEnvironment{dia}{Verbatim}%
  {fontsize=\scriptsize, baselinestretch=0.95}

% \leavevmode mo hop ngang TRUOC khi doi mau.  Thieu no, \color sinh mot
% whatsit o DAU danh sach doc cua o p{} — \vtop khi do cao 0 — nen ca o tut
% xuong dung mot chieu cao dong (~6,6pt) so voi o ben canh.  Day la nguyen
% nhan "lech chu" o moi bang ma cot dau bat dau bang \hl / \no / \yes / \dd.
\newcommand{\dd}[1]{\leavevmode{\color{mute}\scriptsize #1}}
\newcommand{\hl}[1]{\leavevmode{\color{navy}\bfseries #1}}
\newcommand{\no}[1]{\leavevmode{\color{warn}\bfseries #1}}
\newcommand{\yes}[1]{\leavevmode{\color{good}\bfseries #1}}

\title{\bfseries quy gán nguồn cho rag hai tầng}
\subtitle{tập đủ tối tiểu \;$+$\; integrated gradients trên đồ thị}
\date{}
\author{}

\begin{document}

% ════════════════════════════════════════════════════════════════════
\begin{frame}[plain]
  \vfill
  {\color{navy}\Large\bfseries quy gán nguồn cho rag hai tầng}\\[0.9em]
  {\large tập đủ tối tiểu \;$+$\; integrated gradients trên đồ thị}\\[2.2em]
  \begin{tabular}[t]{@{}ll@{}}
    \color{mute}tầng a & tìm \hl{tập đủ tối tiểu} trong pool ứng viên\\
                       & \dd{mwc · deegan–packel · banzhaf có dấu}\\[0.5em]
    \color{mute}tầng b & tìm \hl{đường} trên đồ thị tri thức\\
                       & \dd{integrated gradients trên gate của cạnh}\\
  \end{tabular}
  \vfill
\end{frame}

% ════════════════════════════════════════════════════════════════════
\section{bài toán}

\begin{frame}{bài toán}
  rag trả lời đúng, nhưng không nói được \emph{vì sao}. hai câu hỏi khác nhau:

  \vspace{0.6em}
  \begin{columns}[t]
  \begin{column}{0.48\textwidth}
    \begin{bx}[title={đoạn nào gánh câu trả lời}]
      trong $n$ đoạn đã lấy về, tập nào \emph{thật sự} đủ để trả lời,
      và không đoạn nào thừa?
      \vspace{0.3em}
      \dd{prompt phình theo số đoạn $\Rightarrow$ ablation}
    \end{bx}
  \end{column}
  \begin{column}{0.48\textwidth}
    \begin{bx}[title={vì sao truy hồi tìm ra chúng}]
      trên đồ thị tri thức, \emph{đường} nào dẫn từ thực thể trong câu hỏi
      tới đoạn đó?
      \vspace{0.3em}
      \dd{đồ thị cố định, hàng nghìn cạnh $\Rightarrow$ ig}
    \end{bx}
  \end{column}
  \end{columns}

  \vspace{0.8em}
  \begin{keyb}
    hai câu hỏi cần hai công cụ khác nhau. ghép chúng bằng \hl{một giao diện duy nhất}:
    tầng a trả về tập đích $T$, tầng b nhận $T$. \no{không điểm số nào đi qua.}
  \end{keyb}
\end{frame}

\begin{frame}[fragile]{toàn cảnh}
\begin{dia}
truy vấn q
   |
   |  3 nguồn -> pool thô K <= 60 -> khử trùng lặp -> rerank THEO TỪNG NGUỒN
   v
pool C(q),  n = 40  =  20 bm25  +  15 gfm  +  5 ngẫu nhiên
   |
   |  >>> TẦNG A   tìm tập đủ tối tiểu            (năm bước a1..a5)
   |      v(S)    model đọc S rồi trả lời — khớp đáp án vàng không?      a1
   |      adq(S)  các đoạn trong S có THẬT SỰ nói ra đáp án không?       a2
   |      tau_q   ngưỡng riêng từng truy vấn, chốt từ trần quan sát      a3
   |      g_theta ước lượng v(S) để khỏi gọi llm cho mọi tập con         a4
   |      tìm     đi lên từ nghiệm phủ, rồi chứng minh tối tiểu          a5
   v
T = S*  =  {đoạn 3, đoạn 12}          <- THỨ DUY NHẤT đi qua ranh giới
   |
   |  >>> TẦNG B   tìm đường trên đồ thị          (ba bước b1..b3)
   |      b1 ĐO    ig trên gate của từng (cạnh, tầng)
   |      b2 NỐI   trải phẳng theo tầng, quy hoạch động max-sum
   |      b3 THỬ   suff / nec bằng model thật
   v
Pi_D  =  {đường 1, đường 2}  +  suff(Pi_D)  +  credit từng đoạn
\end{dia}
\end{frame}

% ════════════════════════════════════════════════════════════════════
\section{tầng a}

\begin{frame}{phương pháp tầng a — năm bước}
  \vspace{0.2em}
  \begin{center}
  \footnotesize
  \begin{tabular}[t]{@{}l>{\raggedright\arraybackslash}p{0.235\linewidth}
                        >{\raggedright\arraybackslash}p{0.205\linewidth}
                        >{\raggedright\arraybackslash}p{0.175\linewidth}
                        >{\raggedright\arraybackslash}p{0.175\linewidth}@{}}
    \toprule
    & \textbf{làm gì} & \textbf{bằng công cụ gì} & \textbf{giá mỗi truy vấn}
    & \textbf{cho ra}\\
    \midrule
    \hl{a1} & đo \emph{trả lời đúng} & llm sinh $+$ llm judge
      & 2 lượt mỗi \emph{tập con} & $v(S) \in [0,1]$\\[0.45em]
    \hl{a2} & đo \emph{có bằng chứng} & chẻ câu hỏi $+$ model nli nhỏ
      & 1 lượt, rồi 0 & bảng phủ, $\mathrm{adq}(S)$\\[0.45em]
    \hl{a3} & chốt ngưỡng riêng cho truy vấn & $v(\emptyset)$, $v(R)$, $v(C)$
      & 3 lượt & $\tau_q$, $v_{\mathrm{ref}}$\\[0.45em]
    \hl{a4} & thay $v$ bằng ước lượng rẻ & surrogate $g_\theta$
      & \yes{0 lượt} & $\hat v(S)$\\[0.45em]
    \hl{a5} & tìm, rồi \emph{chứng minh} tối tiểu & tầng phủ $+$ beam $+$ quét dàn
      & ${\sim}6$ lượt & $T = S^{*}$\\
    \bottomrule
  \end{tabular}
  \end{center}

  \vspace{0.5em}
  \begin{keyb}
    hai chỗ quyết định toàn bộ chi phí. \hl{a2} chạy \emph{trước} \hl{a1} dù nó kiểm
    điều kiện thứ hai — vì nó miễn phí và loại ${\sim}45\%$ truy vấn ngay.
    \hl{a4} tồn tại chỉ để $2^{40}$ tập con không phải đi qua llm.
    \dd{\;· mọi slide sau đây mang nhãn bước ở góc phải tiêu đề}
  \end{keyb}
\end{frame}
% ════════════════════════════════════════════════════════════════════
%  tầng a — phần 1:  nền trò chơi · định nghĩa · v · adq · ngưỡng
% ════════════════════════════════════════════════════════════════════

\begin{frame}[fragile]{nền trò chơi hợp tác — bốn khái niệm, một ví dụ}
  \framesubtitle{a · nền}
  ba đoạn $a, b, c$. $v(S)$ là điểm khi cho model đọc đúng tập $S$.

  \vspace{0.4em}
  \begin{columns}[t]
  \begin{column}{0.455\textwidth}
    \begin{bx}[title={ví dụ chạy suốt phần này}]
      \begin{dia}
v(rỗng)      0,05   ngưỡng tau_q = 0,84
v({a})       0,12
v({b})       0,18
v({c})       0,71   <- cao mà KHÔNG có dữ kiện
v({a,b})     0,94   <- bắc cầu
v({a,c})     0,74
v({b,c})     0,76
v({a,b,c})   0,95
      \end{dia}
      \vspace{0.2em}
      tập thắng: $\{a,b\}$, $\{a,b,c\}$. \hl{tối tiểu: chỉ $\{a,b\}$} —
      $\{a,b,c\}$ thắng nhưng chứa một tập thắng nhỏ hơn.
    \end{bx}
  \end{column}
  \begin{column}{0.505\textwidth}
    \begin{bx}[title={bốn khái niệm dùng ngay ở slide sau}]
      \footnotesize
      \begin{tabular}[t]{@{}l@{\hspace{1.1em}}>{\raggedright\arraybackslash}p{0.73\linewidth}@{}}
        \hl{trò chơi} & $(N, v)$ — $N$ là pool, $v$ là hàm giá trị\\[0.45em]
        \hl{đơn điệu} & $S \subseteq S' \Rightarrow v(S) \le v(S')$.
                        \no{ta KHÔNG có tính chất này}\\[0.45em]
        \hl{thắng}    & $v(S) \ge \tau_q$\\[0.45em]
        \hl{mwc}      & thắng, và không tập con thật sự nào thắng\\
      \end{tabular}
    \end{bx}
    \begin{badb}[title={$v$ không đơn điệu, và ta không sửa}]
      ví dụ trên tình cờ đơn điệu; thực tế \no{không} — ta gặp đúng ca đó ở slide
      \hl{tối tiểu} ($v(\{a\}) = 0{,}85$ mà $v(\{a,b\}) = 0{,}60$).
      định nghĩa mwc \yes{không cần đơn điệu}, và bao đơn điệu
      $\bar v(S) = \max_{S'\subseteq S} v(S')$ cho $\mathrm{mwc}(\bar v) = W$ —
      \no{không đổi gì ngoài chi phí.}
    \end{badb}
  \end{column}
  \end{columns}
\end{frame}

\begin{frame}{đối tượng cần tìm}
  \framesubtitle{a · mục tiêu}
  \[
    \boxed{\;
      S \text{ đủ} \;\iff\; \underbrace{v(S) \ge \tau_q}_{\text{trả lời đúng}}
      \;\wedge\; \underbrace{\mathrm{adq}(S)}_{\text{có bằng chứng}}
    \;}
  \]
  \[
    S \text{ tối tiểu} \;\iff\; S \text{ đủ} \;\wedge\;
    \nexists\, S' \subsetneq S \text{ đủ}
    \qquad\qquad
    W = \{\, S \text{ đủ tối tiểu} \,\}
    \qquad\qquad
    T = S^{*} \in W
  \]

  \vspace{0.6em}
  \begin{columns}[t]
  \begin{column}{0.48\textwidth}
    \begin{bx}[title={dư thừa hiện ra thành bội số mwc}]
      ba đoạn trùng ý $\Rightarrow$ ba mwc khác nhau, mỗi tập chứa một đoạn.
      \vspace{0.3em}

      nên bản này \no{không cần gom cụm để hạ} $2^K$ — dư thừa tự lộ trong cấu trúc $W$,
      còn $2^K$ thì surrogate \hl{a4} lo.
    \end{bx}
  \end{column}
  \begin{column}{0.48\textwidth}
    \begin{bx}[title={cái giá: bỏ shapley là bỏ completeness}]
      shapley có $\sum_c \varphi_c = F(E) - F(E')$ — một đẳng thức chính xác, bắt được
      mọi lỗi cài đặt.
      \vspace{0.3em}

      mwc không có đẳng thức đó. đổi lại: được cả \emph{họ} lời giải thích, chỉ số
      cấu trúc, và độ tin cậy. \dd{tầng b lấy lại đẳng thức này bằng $c/h$}
    \end{bx}
  \end{column}
  \end{columns}
\end{frame}

\begin{frame}[fragile]{«đủ» có hai thành phần, và $v$ chỉ đo một}
  \framesubtitle{a1 · a2}
  \begin{columns}[t]
  \begin{column}{0.425\textwidth}
    \begin{badb}[title={ca mồi tham số}]
      \begin{dia}
q  ai đạo diễn phim thắng
   oscar 1994?
y  Robert Zemeckis

đoạn 19  "Forrest Gump là hiện
          tượng văn hóa thập
          niên 1990"
   -> KHÔNG nói đạo diễn
   -> KHÔNG nói thắng giải

model thấy tên phim
   -> TỰ NHỚ -> trả lời ĐÚNG
v({19}) = 0,71   cao
      \end{dia}
      \vspace{0.2em}
      định nghĩa chỉ-dùng-$v$ \no{nhận đoạn 19 vào $T$}, rồi dạy retriever đi lấy nó.
    \end{badb}
  \end{column}
  \begin{column}{0.535\textwidth}
    \begin{bx}[title={và lệch này đã được đo}]
      \footnotesize
      \begin{tabular}[t]{@{}>{\raggedright\arraybackslash}p{0.76\linewidth}r@{}}
        \multicolumn{2}{@{}l@{}}{\hl{model trả lời đúng khi ngữ cảnh không đủ}}\\[0.15em]
        \quad trên các bộ đo trong bài & 35--62\,\%\\[0.6em]
        \multicolumn{2}{@{}l@{}}{\hl{trong các ca model trả lời sai}}\\[0.15em]
        \quad do \emph{không có} ngữ cảnh & 10,2\,\%\\
        \quad do ngữ cảnh \emph{có mà không đủ} & 66,1\,\%\\[0.6em]
        \multicolumn{2}{@{}l@{}}{\hl{sau truy hồi, tỉ lệ ca thật sự đủ ngữ cảnh}}\\[0.15em]
        \quad musique-ans & 46,2\,\%\\
        \quad hotpotqa & 44,6\,\%\\
      \end{tabular}
      \vspace{0.4em}

      \dd{joren et al., iclr 2025 — arXiv:2411.06037}
    \end{bx}
    \begin{keyb}
      hai bộ ta nhắm: hơn \emph{nửa} truy vấn \no{không có tập đủ trong pool.}
      nên điều kiện thứ hai không phải cầu toàn — nó là ca \emph{đa số}.
    \end{keyb}
  \end{column}
  \end{columns}
\end{frame}

\begin{frame}{hai điều kiện, bốn ô — mỗi ô một nguyên nhân khác}
  \framesubtitle{a1 · a2}
  \vspace{0.5em}
  \begin{center}
  \begin{tabular}[t]{@{}>{\raggedright\arraybackslash}p{0.12\linewidth}
                        >{\raggedright\arraybackslash}p{0.37\linewidth}
                        >{\raggedright\arraybackslash}p{0.37\linewidth}@{}}
    \toprule
    & \textbf{adq phủ hết} & \textbf{adq không phủ}\\
    \midrule
    \hl{$v \ge \tau_q$}
      & \yes{nhãn tốt} — nhận
      & \no{mồi tham số} — loại\\[0.7em]
    \hl{$v < \tau_q$}
      & bằng chứng \emph{có} mà model không nối được\newline
        $\rightarrow$ lỗi \textbf{model sinh}, tầng b \no{không sửa được}
      & pool \emph{thiếu}\newline
        $\rightarrow$ thoát sớm ở \texttt{a0}/\texttt{c0}\\
    \bottomrule
  \end{tabular}
  \end{center}

  \vspace{0.9em}
  \begin{keyb}
    giữ $\mathrm{adq}$ \hl{tách khỏi} $v$, không gộp thành một điểm số — gộp thì mất bảng
    này, và mất luôn khả năng tách \emph{lỗi truy hồi} khỏi \emph{lỗi model sinh}.
    ô dưới-trái là ô mà bản chỉ-dùng-$v$ không bao giờ thấy.
  \end{keyb}
\end{frame}

\begin{frame}[fragile]{$v(S)$ — hai lượt llm, một con số}
  \framesubtitle{a1}
  \begin{columns}[t]
  \begin{column}{0.52\textwidth}
    \begin{bx}[title={bước 1 · sinh từ đúng $S$}]
      \begin{dia}
Answer using ONLY the passages below.
If insufficient, say: INSUFFICIENT.

[1] <nội dung chunk thứ nhất trong S>
[2] <nội dung chunk thứ hai>
Question: <q>
Answer:
      \end{dia}
      \vspace{-0.2em}
      \[ \hat a_S = \mathrm{generate}(q, S) \]
    \end{bx}
    \begin{bx}[title={bước 2 · judge, đọc logprob token đầu}]
      \begin{dia}
Question: <q>
Gold answer: Robert Zemeckis
Also acceptable: Zemeckis · Robert Lee
Candidate answer: <â_S>
Correct? One word: Yes or No.
Answer:
      \end{dia}
    \end{bx}
  \end{column}
  \begin{column}{0.44\textwidth}
    \[
      v(S) \;=\;
      \frac{\sum_{t \in \mathcal{Y}} e^{\ell_t}}
           {\sum_{t \in \mathcal{Y}} e^{\ell_t} + \sum_{t \in \mathcal{N}} e^{\ell_t}}
      \;\in [0,1]
    \]
    \dd{$\ell_t$ = logprob token đầu · $\mathcal{Y}$ = \{Yes, yes, ...\}}

    \vspace{0.5em}
    \begin{bx}[title={ba chỗ không được làm sai}]
      \footnotesize
      \no{không} log-likelihood của $y$ — \emph{"Paris"} vs \emph{"Paris, France"} đúng
      ngữ nghĩa mà logprob thấp, và lệch đó \emph{không} triệt tiêu giữa các tập con:
      \textbf{thứ tự lật}.
      \vspace{0.35em}

      \no{không} thang 1--5 — 5 mức làm hàng chục tập con hòa nhau, không phân biệt
      được $S$ với $S \setminus \{c\}$, và \emph{tối tiểu} mất nghĩa.
      \vspace{0.35em}

      \hl{phải} nhét \texttt{answer\_aliases} vào prompt judge, không thì tập \emph{đúng}
      bị loại oan.
    \end{bx}
    \vspace{-0.3em}
    \dd{$\mathrm{cost}(S) = \text{overhead} + \sum_{k \in S} \mathrm{len}(k)$ —
        không phải $500 + 250|S|$, vì document ${\sim}100$ và summary ${\sim}600$
        là hai lớp độ dài}
  \end{column}
  \end{columns}
\end{frame}

\begin{frame}[fragile]{$\mathrm{adq}$ — chẻ câu hỏi, rồi đối chiếu}
  \framesubtitle{a2}
  \begin{columns}[t]
  \begin{column}{0.46\textwidth}
    \begin{bx}[title={bước 1 · chẻ thành mảnh}]
      \footnotesize
      \hl{hỏi ai:} musique gán sẵn phân rã 1-hop $+$ đáp án trung gian.
      corpus riêng thì llm, \textbf{1 lượt mỗi truy vấn} — không phải mỗi tập con.
      \begin{dia}
mảnh 1  phim thắng oscar 1994 là
        Forrest Gump
mảnh 2  Forrest Gump do Zemeckis
        đạo diễn
      \end{dia}
    \end{bx}
    \begin{bx}[title={bước 2 · đối chiếu}]
      \footnotesize
      \hl{hỏi ai:} model nli nhỏ (${\sim}400$M), chạy local, \textbf{0 lượt api}.
      \[ e(k,j) = \mathrm{nli}\bigl(\text{đoạn } k,\; \text{mảnh } j\bigr) \]
      \vspace{-0.5em}

      \dd{đoạn $k$ vào chỗ \emph{premise}, mảnh $j$ vào chỗ \emph{hypothesis} ·
          $40 \times J$ lượt forward model nhỏ}
    \end{bx}
  \end{column}
  \begin{column}{0.5\textwidth}
    \begin{bx}[title={bảng phủ — 0 lượt llm}]
      \begin{dia}
            mảnh 1   mảnh 2
đoạn  3     entail      —
đoạn  7     entail      —
đoạn 12       —      entail
đoạn 23       —      entail
đoạn 31       —      entail
đoạn 19       —         —      <- bẫy
      \end{dia}
      \vspace{-0.3em}
      \begin{align*}
        \mathrm{cov}(S) &= \{\, j : \exists k \in S,\; e(k,j) = \text{entail} \,\}\\
        \mathrm{adq}(S) &\iff \mathrm{cov}(S) = \{1, \dots, J\}
      \end{align*}
    \end{bx}
    \begin{keyb}
      \footnotesize
      bảng phủ cho \textbf{ba} thứ, cả ba miễn phí:
      \begin{dia}
1 CỬA THOÁT  mảnh nào trống -> bỏ ngay
2 ỨNG VIÊN   phủ tập nhỏ, giải CHÍNH XÁC
             {3,12} {3,23} {7,12} ...
3 PRIOR |T|  |T| nên bằng J
      \end{dia}
    \end{keyb}
  \end{column}
  \end{columns}
\end{frame}

\begin{frame}{vì sao phải chẻ ra — hai lý do độc lập}
  \framesubtitle{a2}
  \begin{columns}[t]
  \begin{column}{0.48\textwidth}
    \begin{bx}[title={độ dài}]
      premise $=$ cả $S$ có thể tới $10{,}000$ token, mà model nli chỉ nhận $512$.
      \vspace{0.4em}

      chẻ ra thì mỗi lượt kiểm là \textbf{1 đoạn} (${\sim}100$ token) đối
      \textbf{1 câu ngắn}.
    \end{bx}
  \end{column}
  \begin{column}{0.48\textwidth}
    \begin{badb}[title={bắc cầu}]
      claim ghép \emph{"phim thắng 1994 do Zemeckis đạo diễn"} thì
      \no{không đoạn nào một mình} khẳng định được.
      \vspace{0.4em}

      không chẻ $\Rightarrow$ loại oan cả $\{3,12\}$, tức loại đúng
      \textbf{ca multi-hop} mà cả dự án nhắm.
    \end{badb}
  \end{column}
  \end{columns}

  \vspace{0.7em}
  \begin{keyb}
    chẻ ra thì mỗi mảnh có \emph{đúng một} đoạn khẳng định nó, và ghép lại thành $T$ đầy đủ.
    \dd{\;· cơ chế theo arXiv:2608.00585 — kiểm từng đoạn đối câu hỏi gốc thất bại ở multi-hop}
  \end{keyb}
\end{frame}

\begin{frame}[fragile]{ngưỡng và trần}
  \framesubtitle{a3}
  \vspace{-0.2em}
  \begin{align*}
    \tau_q &= \theta \,\bigl[\, v_{\mathrm{ref}}(q) - v(\emptyset) \,\bigr] + v(\emptyset),
      \qquad \theta = 0{,}95\\[-0.1em]
    v_{\mathrm{ref}}(q) &= \max\Bigl(\, v(C),\; v(R),\; \max_k v(\{k\}),\;
      \text{mọi } v(S) \text{ đã quan sát} \Bigr)
  \end{align*}
  \vspace{-0.6em}
  \begin{columns}[t]
  \begin{column}{0.48\textwidth}
    \begin{bx}[title={mốc phải là cực đại quan sát được}]
      khi $C$ còn đoạn nhiễu thì một tập con có thể ghi điểm \emph{cao hơn} cả $C$ —
      neo vào $v(C)$ là neo vào \no{mốc bị dìm}.
      \begin{dia}
v(C)  = 0,81   nhồi 40 đoạn -> phân tán
v(R)  = 0,88   12 đoạn, ít nhiễu hơn
->  v_ref = 0,88,  tau_q = 0,84
      \end{dia}
    \end{bx}
  \end{column}
  \begin{column}{0.48\textwidth}
    \begin{bx}[title={hai mốc đặc biệt}]
      \begin{dia}
v(rỗng)  không cho đoạn nào
         -> chỉ tri thức tham số
v(R)     12 đoạn: top-2 mỗi cụm, VÀ
         tôn trọng hạn ngạch 20/15/5
      \end{dia}
      \vspace{0.2em}
      \dd{\textbf{chốt chặn mẫu số nhỏ:}\;
          $v_{\mathrm{ref}} - v(\emptyset) < 0{,}05 \Rightarrow T = \emptyset$,
          cờ \emph{"truy hồi không đóng góp"} — ca có thật, không phải lỗi.}
    \end{bx}
  \end{column}
  \end{columns}
\end{frame}

\begin{frame}{cập nhật $v_{\mathrm{ref}}$ tự nhất quán — và $\theta < 1$ gánh tất}
  \framesubtitle{a3}
  quan sát mới có thể \emph{cao hơn} trần đang dùng. lúc đó $v_{\mathrm{ref}}$ tăng,
  nên $\tau_q$ cũng tăng — liệu tập vừa nâng trần có bị chính ngưỡng mới loại?

  \vspace{0.8em}
  \[
    v(S) > v_{\mathrm{ref}}
    \;\Longrightarrow\;
    v_{\mathrm{ref}} := v(S)
    \;\Longrightarrow\;
    \tau_q' = \theta\bigl[v(S) - v(\emptyset)\bigr] + v(\emptyset)
    \;\underbrace{<}_{\theta\,<\,1}\; v(S)
  \]

  \vspace{0.8em}
  \begin{columns}[t]
  \begin{column}{0.48\textwidth}
    \begin{keyb}[title={không bao giờ tự phủ định}]
      tập vừa nâng trần \yes{luôn vẫn vượt ngưỡng mới}. cập nhật không thể loại
      chính cái gây ra nó.
    \end{keyb}
  \end{column}
  \begin{column}{0.48\textwidth}
    \begin{keyb}[title={không phải xét lại lịch sử}]
      $\tau_q$ chỉ \emph{tăng}, nên mọi tập đã bị loại vẫn bị loại.
      \vspace{0.25em}

      \dd{$\theta = 1$ thì suy biến: $\tau_q' = v(S)$, tập nâng trần chỉ \emph{bằng}
          ngưỡng, và tính chất trên mất}
    \end{keyb}
  \end{column}
  \end{columns}
\end{frame}
% ════════════════════════════════════════════════════════════════════
%  tầng a — phần 2:  surrogate · tìm · tối tiểu · chỉ số
%  rồi mở sang tầng b
% ════════════════════════════════════════════════════════════════════

\begin{frame}[fragile]{surrogate $g_\theta$ — một hàm, hai pha}
  \framesubtitle{a4}
  \[
    g_\theta\Bigl(q,\; S,\; \{\hat s(q,k)\}_{k \in S},\; \{\tilde v_k\}_{k \in S},\; v(\emptyset)\Bigr)
    \;\longrightarrow\; \hat v(S)
  \]
  \dd{bất biến hoán vị trên $S$ · $\hat s$ = điểm rerank, luôn có ·
      $\tilde v_k = v(\{k\}) - v(\emptyset)$, chỉ khi đã trả vòng 0}

  \vspace{0.5em}
  \begin{columns}[t]
  \begin{column}{0.47\textwidth}
    \begin{badb}[title={$\tau_q$ không được là đầu vào}]
      $g_\theta$ hồi quy lên $v(S)$; so với ngưỡng xảy ra \emph{sau}, bên ngoài mô hình.
      \vspace{0.35em}

      bước \texttt{c2} gọi $g_\theta$ để quyết định có mua $v(C)$ — mà $v(C)$ mới chốt
      $v_{\mathrm{ref}}$ rồi chốt $\tau_q$. cho $\tau_q$ vào đầu vào là \no{tạo vòng tròn thật}.
      \vspace{0.35em}

      $\tau_q$ xuất hiện trong \emph{mất mát}, không ở đầu vào:
      \[ \mathcal{L} = \mathcal{L}_{\text{cls}} + \lambda \mathcal{L}_{\text{rank}} \]
      \dd{hinge theo cặp \emph{trong một} truy vấn}
    \end{badb}
  \end{column}
  \begin{column}{0.49\textwidth}
    \begin{bx}[title={dữ liệu — là log của chính các lượt gọi $v$}]
      \begin{dia}
nguồn                 truy vấn  nhãn      tổng
pha 1, bốn vòng         10.000  x111  1.110.000
verification pha 2     188.000  x 10  1.880.000
tỉa, |S*| = 6..10        8.000  x 13    104.000
cứu -> pipeline đầy đủ   7.400  x111    821.000
trượt bị bỏ             14.000  x  5     70.000
                                     ---------
                                     3.983.000
      \end{dia}
      \vspace{0.2em}
      $|S|$ chạy từ 0 tới 40 — \hl{không cỡ nào đặc biệt}, nên \emph{một} mô hình,
      không phải một cho singleton và một cho phần còn lại.
    \end{bx}
  \end{column}
  \end{columns}
\end{frame}

\begin{frame}[fragile]{che $\tilde v_k$ theo từng chunk, và nghiệm thu}
  \framesubtitle{a4}
  \begin{columns}[t]
  \begin{column}{0.49\textwidth}
    \begin{bx}[title={ba khuôn, vì thực tế có ba chế độ}]
      \begin{dia}
pha 1 + cứu        cả 40 chunk có v~_k
verification pha 2 CHỈ |T| chunk có
tìm theo tầng (c5) không chunk nào có
      \end{dia}
      \vspace{0.3em}
      trộn ba khuôn theo đúng tỉ lệ chúng xuất hiện.
      \vspace{0.3em}

      che tất-cả-hoặc-không-gì \no{không dạy được khuôn giữa} — mà đó là khuôn của
      \emph{mọi} lượt verification ở pha 2.
    \end{bx}
  \end{column}
  \begin{column}{0.47\textwidth}
    \begin{keyb}[title={nghiệm thu, và sàn phải vượt}]
      \hl{auc trong từng truy vấn}, rồi trung bình qua các truy vấn.
      \vspace{0.3em}

      \hl{sàn:} mô hình chỉ nhìn $|S|$.
      \vspace{0.3em}

      auc \emph{gộp} là bẫy: tập lớn thì $v$ cao nên con số đẹp sẵn, mà phân biệt
      \emph{trong} truy vấn bằng \no{0} — mọi tập cỡ 3 hòa điểm, beam thành chọn bừa.
    \end{keyb}
    \begin{bx}[title={mốc kiến trúc phải vượt}]
      gbm trên đặc trưng tay chạy \hl{trước}. set transformer chỉ đáng chọn nếu
      vượt nó trên auc-trong-truy-vấn.
    \end{bx}
  \end{column}
  \end{columns}
\end{frame}

\begin{frame}[fragile]{tìm theo tầng, đi lên từ nghiệm phủ}
  \framesubtitle{a5}
  \begin{columns}[t]
  \begin{column}{0.45\textwidth}
    \begin{keyb}[title={vì sao chỉ tìm trong vùng này}]
      \begin{dia}
tập KHÔNG phủ hết mảnh
  -> nếu vượt tau_q thì chỉ có thể
     vì mồi tham số
  -> mà ta đã quyết LOẠI ca đó
  -> tìm ở đó là tìm thứ ta sẽ bỏ
      \end{dia}
      \vspace{0.2em}
      cả dàn $2^n$ co lại thành \hl{nghiệm phủ, và mọi tập chứa một nghiệm phủ}.
    \end{keyb}
    \begin{bx}[title={dự bị}]
      bảng phủ có thể thiếu cột (nli bỏ sót bằng chứng diễn đạt lại). mọi tầng
      trượt $\Rightarrow$ quay về beam mò trên toàn dàn, tỉa theo $r_i$ tăng dần.
    \end{bx}
  \end{column}
  \begin{column}{0.51\textwidth}
    \begin{bx}[title={các tầng}]
      \begin{dia}
tầng 0  nghiệm phủ tối tiểu       |T| = J
        giải CHÍNH XÁC từ bảng phủ
tầng 1  nghiệm phủ + 1 đoạn
tầng 2  nghiệm phủ + 2 đoạn
...
mỗi tầng: chấm bằng g_theta, lấy top-beam,
          đo v theo thứ tự đó
      \end{dia}
    \end{bx}
    \begin{bx}[title={hai lợi ích rơi ra từ cấu trúc}]
      \textbf{1}\; mọi ứng viên đã đủ về bằng chứng — không còn ứng viên mồi tham số.
      \vspace{0.25em}

      \textbf{2}\; thứ tự theo \emph{cỡ}, nên tập nhỏ luôn thử trước:
      \yes{"tối tiểu trước" là miễn phí theo cấu trúc}, không phải một bước tỉa riêng.
    \end{bx}
  \end{column}
  \end{columns}
\end{frame}

\begin{frame}[fragile]{tối tiểu — vì sao phải quét hết, và vì sao không tốn}
  \framesubtitle{a5}
  \begin{columns}[t]
  \begin{column}{0.47\textwidth}
    \begin{badb}[title={bỏ-một-phần-tử là SAI, vì $v$ không đơn điệu}]
      \begin{dia}
tau_q = 0,76
v({a})     = 0,85  >= tau_q  <- {a} ĐÃ ĐỦ
v({a,b})   = 0,60  <  tau_q     b làm lạc
v({a,c})   = 0,71  <  tau_q
v({b,c})   = 0,40  <  tau_q
v({a,b,c}) = 0,80  >= tau_q

bỏ-một-phần-tử:  ba cặp đều trượt
                 -> tuyên "tối tiểu"  SAI
sự thật:         {a} thắng -> thừa b, c
      \end{dia}
    \end{badb}
  \end{column}
  \begin{column}{0.49\textwidth}
    \begin{keyb}[title={nhưng chỉ tập con CÒN PHỦ mới đe dọa}]
      \begin{dia}
T = {3, 12, 19}
  3 phủ mảnh 1 · 12 phủ mảnh 2 · 19 thừa

tập con thực sự:        7 tập
trong đó CÒN phủ hết:   1 tập -- {3,12}
  {3,19} mất mảnh 2 · {12,19} mất mảnh 1
  mọi đơn tử mất một mảnh
      \end{dia}
      \vspace{0.2em}
      và tầng 0 \yes{đã đo $\{3,12\}$ rồi} — nên tính tối tiểu \hl{không tốn lượt nào thêm}.
    \end{keyb}
    \begin{bx}[title={khi không có bảng phủ: $2^{|S^*|} - 2$ lượt}]
      \footnotesize
      \setlength{\tabcolsep}{3pt}
      \begin{tabular}[t]{@{}lrrrrr@{}}
        \toprule
        $|S^*|$ & 3 & 5 & 6 & 8 & 10\\
        \midrule
        quét hết & \$0,02 & \$0,11 & \$0,24 & \$1,14 & \$5,37\\
        tỉa rồi quét & — & \$0,11 & \$0,04 & \$0,07 & \$0,10\\
        \bottomrule
      \end{tabular}
    \end{bx}
  \end{column}
  \end{columns}
\end{frame}

\begin{frame}[fragile]{ba chỉ số trò chơi — tính thẳng trên ví dụ $a, b, c$}
  \framesubtitle{a5 · nền}
  \[
    \beta_i^{+} = \bigl|\{\, S \subseteq N \setminus \{i\} :
      v(S) < \tau_q \;\wedge\; v(S \cup \{i\}) \ge \tau_q \,\}\bigr|
    \qquad
    \beta_i = \frac{\beta_i^{+} - \beta_i^{-}}{2^{\,n-1}}
    \qquad
    \mathrm{pgi}_i = \bigl|\{\, S \in W : i \in S \,\}\bigr|
  \]
  \dd{$\beta_i^{+}$: thêm $i$ biến thiếu $\to$ đủ \;·\; $\beta_i^{-}$: biến đủ $\to$ thiếu,
      $=0$ nếu đơn điệu \;·\; pgi chỉ đếm trên các mwc}

  \vspace{0.2em}
  \begin{columns}[t]
  \begin{column}{0.53\textwidth}
    \begin{bx}[title={đếm tay, $\tau_q = 0{,}84$, $n = 3$}]
      \begin{dia}
i = a  {b}   0,18 -> {a,b} 0,94  pivotal
       {b,c} 0,76 -> 0,95        pivotal
       hai ca kia: không           B+ = 2
i = b  {a}   0,12 -> {a,b} 0,94  pivotal
       {a,c} 0,74 -> 0,95        pivotal
       hai ca kia: không           B+ = 2
i = c  {a,b} đã thắng -> KHÔNG tính
       ba ca kia: vẫn thua         B+ = 0
      \end{dia}
    \end{bx}
  \end{column}
  \begin{column}{0.43\textwidth}
    \begin{keyb}[title={kết quả}]
      \footnotesize
      \begin{tabular}[t]{@{}lrrr@{}}
        \toprule
        & $a$ & $b$ & $c$\\
        \midrule
        $v(\{i\})$     & 0,12 & 0,18 & \no{0,71}\\
        $\beta_i$      & 0,50 & 0,50 & \no{0}\\
        $\mathrm{pgi}_i$ & 1 & 1 & \no{0}\\
        \bottomrule
      \end{tabular}
    \end{keyb}
    \begin{badb}[title={đây chính là điều cần thấy}]
      $c$ có \emph{điểm đơn lẻ cao nhất} ($0{,}71$) mà \no{cả hai chỉ số đều cho 0} —
      vì $c$ \yes{không bao giờ} là thứ biến một tập thiếu thành đủ.
    \end{badb}
  \end{column}
  \end{columns}
\end{frame}

\begin{frame}[fragile]{chỉ số cấu trúc, và đầu ra}
  \framesubtitle{a5}
  \begin{columns}[t]
  \begin{column}{0.47\textwidth}
    % KHONG dat \dd (hay bat ky van ban tieng Viet nao) ben trong align*/$...$:
    % no chay o font TOAN, va cmmi10 khong co glyph co dau — chu bi NUOT im
    % lang, bien dich van sach.  "trung lap giua cac mwc" in ra "trnglpgiaccmwc".
    \[ \mathrm{ov}(S) = \sum_{S' \ne S} \frac{|S \cap S'|}{|S \cup S'|} \]
    \vspace{-1.1em}

    \dd{độ trùng lặp giữa các mwc — cao nghĩa là nhiều lời giải thích gần như nhau}
    \vspace{0.05em}
    \begin{bx}[title={$\beta_i < 0$ chỉ thẳng ra đoạn gây hại}]
      đúng vai mà $\varphi_c < 0$ đóng ở bản shapley: \no{không đơn điệu} ở đây là
      \yes{thông tin}, không phải khuyết tật.
      \vspace{0.2em}

      \hl{chuẩn hóa $\beta$ và pgi riêng rồi mới gộp} — $\beta \in [-1,1]$ còn pgi
      đếm tới hàng chục, cộng thẳng thì pgi nuốt $\beta$.
    \end{bx}
    \begin{bx}[title={chọn $S^{*}$}]
      \begin{dia}
parsimony   -|S|
structure   sigma_g = w_b*beta + w_p*PGI
rank        tổng r_i   (chỉ phá hòa cuối)
-> lọc Pareto, rồi gộp có trọng số
      \end{dia}
    \end{bx}
  \end{column}
  \begin{column}{0.49\textwidth}
    \begin{keyb}[title={đầu ra}]
      \begin{dia}
T = S*             <- THỨ DUY NHẤT qua ranh giới
+ phiên bản g_theta   lô nào gán
+ đường đi            qua · tỉa · thử lại · cứu
+ trọng số            N_s / n_s
+ v(T) - tau_q        khoảng vượt ngưỡng
+ cov(T), J           phủ mấy trên mấy mảnh

W, p(S), entropy(p)  báo cáo, độ tin cậy
beta_i < 0           ứng viên V_neg (tầng b)
      \end{dia}
      \vspace{0.2em}
      \dd{entropy thấp = một người thắng rõ · cao = nhiều mwc suýt soát}
    \end{keyb}
  \end{column}
  \end{columns}
\end{frame}

% ════════════════════════════════════════════════════════════════════
\section{tầng b}

\begin{frame}[fragile]{bài toán tầng b — thứ cần đọc ra là một CHUỖI}
  \framesubtitle{b · mục tiêu}
  \begin{columns}[t]
  \begin{column}{0.46\textwidth}
    $T$ từ tầng a, và đồ thị. gfm truyền tin $L$ bước từ node khởi tới $T$.
    \hl{câu hỏi:} \emph{đường} nào dẫn tới $D \in T$?
    \vspace{0.5em}
    \begin{badb}[title={vì sao không chỉ chấm điểm từng cạnh}]
      điểm từng cạnh cho một \emph{đống} cạnh quan trọng, không cho một \emph{chuỗi}.
      \vspace{0.3em}

      mà lời giải thích đọc được phải là
      \begin{dia}
q -> Forrest Gump -> Zemeckis -> D
      \end{dia}
      tức một \hl{dãy cạnh nối tiếp}, mỗi bước ở đúng tầng của nó.
    \end{badb}
  \end{column}
  \begin{column}{0.50\textwidth}
    \begin{keyb}[title={ba bước — xương sống của cả phần này}]
      \begin{dia}
b1 ĐO   ig trên gate của từng (cạnh, tầng)
        -> cạnh nào gánh, ở bước nào

b2 NỐI  trải phẳng đồ thị theo tầng rồi
        quy hoạch động max-sum
        -> dãy cạnh nối tiếp, không phải đống

b3 THỬ  suff / nec bằng model THẬT
        -> đường tìm được có thật đủ không
      \end{dia}
    \end{keyb}
    \begin{bx}[title={điểm mục tiêu}]
      \[ F_B \;=\; \log p(v_D) \]
      \dd{$p(v_D)$: xác suất gfm gán cho node đích $D$ · một lượt forward cho ra
          $p(v)$ của \emph{mọi} node, nên $F_B$ của cả $T$ dùng chung một lượt}
    \end{bx}
  \end{column}
  \end{columns}
\end{frame}

\begin{frame}[fragile]{integrated gradients — công cụ của bước $\mathtt{b1}$}
  \framesubtitle{b1 · nền}
  gradient tại một điểm chỉ nói độ dốc \emph{ở đó}. ig tích phân dọc đường từ mốc tới điểm thật:

  \[
    \mathrm{ig}_i \;=\; (x_i - x_i^{0}) \times
    \frac{1}{M}\sum_{m=1}^{M}
    \frac{\partial F\!\left(x^{0} + \alpha_m (x - x^{0})\right)}{\partial x_i},
    \qquad \alpha_m = \tfrac{m}{M}
  \]

  \vspace{0.3em}
  \begin{columns}[t]
  \begin{column}{0.48\textwidth}
    \begin{bx}[title={tính chất ta dùng}]
      \textbf{completeness}\;
      $\sum_i \mathrm{ig}_i = F(x) - F(x^0)$
      \vspace{0.3em}

      tổng đóng góp \emph{bằng đúng} phần mà đầu vào tạo ra. nên credit chia được, và
      kiểm được bằng một đẳng thức.
      \vspace{0.3em}

      \dd{đây là đẳng thức mà tầng a đã mất khi bỏ shapley — tầng b lấy lại được}
    \end{bx}
  \end{column}
  \begin{column}{0.48\textwidth}
    \begin{bx}[title={vì sao ig cho cạnh, ablation cho đoạn}]
      \[ \text{ig rẻ hơn ablation} \iff N > 3M \]
      \vspace{-0.6em}
      \begin{dia}
cạnh đồ thị  hàng nghìn, đồ thị CỐ ĐỊNH -> ig
token        hàng trăm, prompt CỐ ĐỊNH  -> ig
đoạn / pool  prompt PHÌNH theo số đoạn
             -> tỉ lệ = 3M, K triệt tiêu
             -> KHÔNG có điểm hòa vốn
             -> ablation
      \end{dia}
    \end{bx}
  \end{column}
  \end{columns}
\end{frame}

\begin{frame}[fragile]{gate trên cạnh — và phải riêng từng tầng}
  \framesubtitle{b1}
  \[
    \text{truyền tin bước } l:\qquad
    h_v^{(l)} \;=\; \mathrm{agg}\Bigl(
      \bigl\{\, g_e^{(l)} \cdot m_e^{(l)} \;:\; e = (u \to v) \,\bigr\}
    \Bigr)
  \]
  \dd{$g_e^{(l)} \in [0,1]$ là \emph{gate} ta chèn vào cạnh $e$ ở tầng $l$ ·
      $g \equiv 1$ cho đúng mô hình gốc}

  \vspace{0.5em}
  \begin{columns}[t]
  \begin{column}{0.48\textwidth}
    \begin{badb}[title={một gate cho mỗi cạnh là KHÔNG đủ}]
      cùng một cạnh có thể cần ở bước 1 mà không cần ở bước 3. gate dùng chung
      \no{trộn hai vai} và ig trả về một số vô nghĩa.
      \vspace{0.3em}

      \begin{dia}
gate chung       g_e        |E| biến
gate theo tầng   g_e^(l)    |E| x L biến
      \end{dia}
      \vspace{0.2em}
      chỉ gate theo tầng mới đọc ra được \emph{thứ tự} các bước — tức mới ra được
      một \hl{chuỗi}.
    \end{badb}
  \end{column}
  \begin{column}{0.48\textwidth}
    \begin{bx}[title={ig trên gate}]
      \[
        \mathrm{ig}_e^{(l)} = \bigl(g^{*(l)}_e - 0\bigr) \times
        \frac{1}{M}\sum_{m=1}^{M}
        \frac{\partial F_B\bigl(\alpha_m g^{*}\bigr)}{\partial g_e^{(l)}}
      \]
      \vspace{-0.4em}
      \begin{dia}
g* = 1    khi ĐÀO NHÃN
g* = ĝ    khi BÁO CÁO
      \end{dia}
      \vspace{0.2em}
      \footnotesize
      hệ số $\hat g$ cắt cả hai chiều: tốt cho báo cáo trung thực, \no{tệ cho đào nhãn} —
      $\hat g$ là thứ model \emph{làm}, không phải thứ ta dùng để \emph{đo}.
      nên lúc mining luôn $g^{*} = 1$.
    \end{bx}
  \end{column}
  \end{columns}
\end{frame}

\begin{frame}[fragile]{tính chất chia đều — kiểm tra nội bộ miễn phí}
  \framesubtitle{b1}
  một đường $\pi$ dùng $h$ cạnh, và đường đó đóng góp $c$ vào $F_B$. completeness của ig cho:
  \[
    \sum_{e \in \pi} \mathrm{ig}_e \;=\; c
    \qquad\text{và nếu các cạnh đối xứng trong } \pi \text{ thì}\qquad
    \mathrm{ig}_e \;=\; \frac{c}{h}
  \]

  \vspace{0.5em}
  \begin{columns}[t]
  \begin{column}{0.48\textwidth}
    \begin{keyb}[title={vì sao đáng tiền}]
      đây là một \hl{đẳng thức kiểm được}, không phải một trực giác.
      \vspace{0.3em}

      cài sai chỗ nào trong khâu ig — sai dấu, sai số bước $M$, sai mốc, nhân
      hai lần hệ số — thì $\sum_{e \in \pi} \mathrm{ig}_e \ne c$ và ta thấy ngay.
      \vspace{0.3em}

      \dd{đúng vai mà $\sum_c \varphi_c = F(E) - F(E')$ đóng ở bản shapley —
          phần bất biến mà tầng a đã mất khi bỏ completeness}
    \end{keyb}
  \end{column}
  \begin{column}{0.48\textwidth}
    \begin{bx}[title={ca một đường, ba cạnh}]
      \begin{dia}
c = 0,81          đóng góp của cả đường

ig cạnh 1 = 0,27
ig cạnh 2 = 0,27      mỗi cạnh = c/h
ig cạnh 3 = 0,27           = 0,81/3
            -----
tổng      = 0,81  = c        KHỚP
      \end{dia}
      \vspace{0.2em}
      lệch $\Rightarrow$ có lỗi cài đặt, \no{dừng, đừng đi tiếp}.
    \end{bx}
  \end{column}
  \end{columns}
\end{frame}
\begin{frame}[fragile]{đồ thị trải phẳng, và quy hoạch động}
  \framesubtitle{b2}
  \begin{columns}[t]
  \begin{column}{0.485\textwidth}
    \begin{bx}[title={vẽ lại: mỗi tầng một bản sao}]
      \begin{dia}
node (v, l)   = node v ở bước l
cạnh (u,l) -> (v,l+1)  nếu u->v trong G

L tầng  ->  đồ thị KHÔNG CÓ CHU TRÌNH
        ->  quy hoạch động một lượt
      \end{dia}
      \vspace{0.2em}
      đồ thị gốc có chu trình, nên không duyệt một lượt được. trải theo tầng thì
      \yes{thứ tự topo là chính $l$}.
    \end{bx}
    \begin{bx}[title={một lượt quét}]
      \[ \mathrm{best}(v, l) = \max_{u \,:\, (u,v) \in E}
         \Bigl[\, \mathrm{best}(u, l-1) + \mathrm{ig}_{(u,v)}^{(l)} \,\Bigr] \]
      \dd{lưu \texttt{back[}$v,l$\texttt{]} để truy ngược ra dãy cạnh ·
          $O\bigl((|V| + |E|)\,L\bigr)$}
    \end{bx}
  \end{column}
  \begin{column}{0.475\textwidth}
    \begin{badb}[title={max-SUM, không phải min-cost}]
      quen tay thì đặt $\mathrm{cost} = -\log(\cdot)$ rồi chạy dijkstra. \no{sai ở đây.}
      \vspace{0.3em}

      \begin{dia}
ig CỘNG TÍNH theo completeness:
   tổng ig trên một đường = đóng góp
   của đường đó

-log thì BIẾN TỔNG thành TÍCH
   -> mất completeness
   -> c/h không còn đúng
   -> mất luôn phép kiểm nội bộ
      \end{dia}
      \vspace{0.2em}
      và ig \emph{có dấu}: cạnh gây hại cho $\mathrm{ig} < 0$, mà $-\log$ của số âm
      thì vô nghĩa.
    \end{badb}
  \end{column}
  \end{columns}
\end{frame}

\begin{frame}[fragile]{kiểm chứng — và hai đại lượng khác nhau}
  \framesubtitle{b2 · b3}
  \[
    \mathrm{suff}(\Pi_D) \;=\;
    \frac{\sum_{(e,l) \,\in\, \Pi_D} \mathrm{ig}_e^{(l)}}
         {\sum_{(e,l)} \mathrm{ig}_e^{(l)}}
    \qquad\qquad
    \mathrm{nec}(\pi) \;=\; F_B(g^{*}) - F_B\bigl(g^{*} \text{ với } \pi \text{ bị xóa}\bigr)
  \]
  \dd{$\mathrm{suff}$: phần của «cái bánh» mà tập đường đã lấy · $\mathrm{nec}$: mất bao nhiêu
      nếu bỏ đúng đường đó · $\mathrm{suff}$ còn là \emph{điều kiện dừng} của bước \texttt{b2}}

  \vspace{0.6em}
  \begin{columns}[t]
  \begin{column}{0.47\textwidth}
    \begin{keyb}[title={$\mathrm{suff}$ tính trên TOÀN tập đường}]
      \no{không} cộng $\mathrm{suff}$ của từng đường rồi lấy tổng — hai đường dùng chung
      cạnh thì cạnh đó bị đếm hai lần.
      \vspace{0.3em}

      lấy \hl{hợp các cạnh} rồi mới chia:
      \begin{dia}
hợp cạnh: 0,22 + 0,19 + 0,33 = 0,74
          0,74 / 0,81 = 91%
suff = 0,91     đủ
      \end{dia}
    \end{keyb}
  \end{column}
  \begin{column}{0.49\textwidth}
    \begin{bx}[title={$\mathrm{suff}$ có thể bị chặn dưới 1 — credit mồ côi}]
      node \emph{không} tới được từ bất kỳ seed nào vẫn góp vào \hl{mẫu số}.
      nên $\mathrm{suff}$ thấp có \emph{hai} nguyên nhân:
      \vspace{0.25em}

      \begin{dia}
suff_max  tính MIỄN PHÍ từ bảng qhđ
          = trần mà mọi đường từ seed
            đạt được

suff_max < theta_s  -> entity linking
                       hỏng, dừng ngay
suff_max >= theta_s -> đồ thị thiếu cạnh
                       (ô «lỗi đồ thị»)
      \end{dia}
    \end{bx}
  \end{column}
  \end{columns}
\end{frame}

\begin{frame}[fragile]{nhiều seed hội tụ — greedy xóa cạnh, rồi tỉa ngược}
  \framesubtitle{b2}
  \begin{columns}[t]
  \begin{column}{0.52\textwidth}
    \begin{bx}[title={hai chặng}]
      \begin{dia}
-- chặng 1: gom ------------------------
Pi_D <- rỗng
lặp:
    pi   = truy ngược back[] từ (D, L)
    gain = tổng ig các cạnh MỚI trong pi
    Pi_D <- Pi_D + pi

    suff(Pi_D) >= theta_s -> DỪNG: thành công
    gain < eps            -> DỪNG: bằng chứng
                             phân tán
    |Pi_D| = n_max        -> DỪNG: hết ngân
                             sách đọc hiểu

    đặt ig = 0 cho mọi cạnh đã lấy
    chạy lại qhđ

-- chặng 2: tỉa ngược ------------------
cho mỗi pi trong Pi_D (theo thứ tự thêm):
    suff(Pi_D \ {pi}) >= theta_s
        -> Pi_D <- Pi_D \ {pi}  đường thừa
      \end{dia}
    \end{bx}
  \end{column}
  \begin{column}{0.44\textwidth}
    \begin{keyb}[title={cạnh dùng chung tự thành miễn phí}]
      đặt $\mathrm{ig} = 0$ cho cạnh đã lấy, nên vòng sau \yes{không trả tiền lại}
      cho cạnh đó. hội tụ nhiều seed xử lý \emph{tự nhiên}, không cần luật riêng.
    \end{keyb}
    \begin{bx}[title={$n_{\max}$ không phải hằng số chọn tay}]
      \[ n_{\max} = \Bigl\lceil \frac{\theta_s}{\tilde\varepsilon} \Bigr\rceil \]
      \dd{$\tilde\varepsilon$: ngưỡng gain tương đối. mỗi đường được nhận phải góp
          $\ge \tilde\varepsilon$, nên quá $\theta_s/\tilde\varepsilon$ đường là
          không thể còn đủ}
      \vspace{0.3em}

      \no{không} phải $|V_q| + 1$ — số seed không bằng số đường độc lập.
    \end{bx}
  \end{column}
  \end{columns}
\end{frame}

% ════════════════════════════════════════════════════════════════════
\section{nối hai tầng}

\begin{frame}[fragile]{chỉ một thứ đi qua ranh giới}
  \begin{columns}[t]
  \begin{column}{0.415\textwidth}
    \begin{keyb}[title={giao diện}]
      \begin{dia}
TẦNG A  ->  T = {D_2, D_7, D_15}  ->  TẦNG B

KHÔNG qua: v(S) · tau_q · beta_i · PGI
           g_theta · W · entropy
      \end{dia}
      \vspace{0.25em}
      điểm của tầng a đo trên \emph{pool}; tầng b định tuyến trên \emph{đồ thị} — hai
      không gian khác nhau. và nhãn $(q \to T)$ \hl{không phụ thuộc kiến trúc gfm}:
      đào một lần, dùng cho mọi lần train.
    \end{keyb}
  \end{column}
  \begin{column}{0.545\textwidth}
    \begin{bx}[title={ba hàm mất mát}]
      \begin{align*}
        \mathcal{L}_{\text{node}} &= \mathrm{bce}\bigl(p(v),\; \mathbf{1}[v \in T]\bigr)\\[0.2em]
        \mathcal{L}_{\text{hop}}  &= \sum_{l}\;\sum_{v \,\in\, V_{\text{hop}}(l)}
          \mathrm{bce}\bigl(p^{(l)}(v),\; \mathbf{1}[(v,l) \in \Pi]\bigr)\\[0.2em]
        \mathcal{L}_{\text{edge}} &= \sum_{(e,l) \,\in\, \Pi}
          -\log \hat g_e^{(l)}
      \end{align*}
      \vspace{-0.3em}
      \begin{dia}
V_hop(l) = { v : (v,l) thuộc Pi }   dương
         + { v : (v,l') thuộc Pi }  âm KHÓ
                 đúng node, SAI tầng
         + V_neg                    âm thường
      \end{dia}
      \dd{không có $V_{\text{hop}}(l)$ thì tỉ lệ dương:âm thành \no{$1 : 10^6$}}
    \end{bx}
  \end{column}
  \end{columns}
\end{frame}

\begin{frame}[fragile]{bảng chẩn đoán — và nó cũng là nguồn nhãn}
  \[
    k_{\text{top}}(q) = \bigl\lceil \beta \,\bigl|T(q)\bigr| \bigr\rceil,
    \quad \beta \approx 1{,}5
    \qquad\qquad
    \mathrm{rank}_{\text{gfm}}(D) \;\text{đọc từ \emph{một} lượt forward}
  \]
  \dd{ngưỡng hạng theo \emph{từng ví dụ}, không phải một $k$ cố định toàn cục}

  \vspace{0.7em}
  \begin{center}
  \begin{tabular}[t]{@{}>{\raggedright\arraybackslash}p{0.11\linewidth}
                        >{\raggedright\arraybackslash}p{0.355\linewidth}
                        >{\raggedright\arraybackslash}p{0.355\linewidth}@{}}
    \toprule
    & \textbf{$\mathrm{rank}$ tốt} & \textbf{$\mathrm{rank}$ kém}\\
    \midrule
    \hl{$\mathrm{suff}$ cao}
      & \yes{lành} — truy hồi tìm đúng, và đường giải thích được
      & \no{lỗi model} — đường có đó mà gfm không xếp $D$ lên cao\\[0.6em]
    \hl{$\mathrm{suff}$ thấp}
      & \no{đường tắt} — xếp đúng \emph{nhờ} lối khác, không nhờ bằng chứng
      & \no{lỗi đồ thị} — thiếu cạnh, hoặc entity linking hỏng\\
    \bottomrule
  \end{tabular}
  \end{center}

  \vspace{0.7em}
  \begin{keyb}
    $\mathrm{suff}$ kiểm \hl{đồ thị}, $\mathrm{rank}$ kiểm \hl{model} — hai trục độc lập.
    ô \emph{đường tắt} là ô dễ bỏ sót nhất: chỉ số đẹp mà lời giải thích rỗng.
    và mỗi ô \yes{tự nó là một lô nhãn} để train vòng sau.
  \end{keyb}
\end{frame}
% ════════════════════════════════════════════════════════════════════
\section{thi hành}

\begin{frame}[fragile]{hai pha — mô hình ở pha 1, nhãn ở pha 2}
  \framesubtitle{thi hành}
  chạy pipeline đầy đủ cho $200{,}000$ truy vấn thì không trả nổi. nên tách hai pha
  có \emph{mục đích} khác nhau:

  \vspace{0.5em}
  \begin{columns}[t]
  \begin{column}{0.46\textwidth}
    \begin{bx}[title={pha 1 · $N_1 = 10{,}000$ · llm đầy đủ}]
      \begin{dia}
A0  chẻ câu hỏi + bảng phủ       400
A1  vòng 0  v(rỗng)·v(R)·v(C)
            + 40 singleton    44.500
A2  vòng 1  greedy biên       27.000
A3  vòng 2  active learning  ~37.000
A4  vòng 3  xác nhận         ~15.000
                            ---------
                            ~123.900 token
      \end{dia}
      \vspace{0.2em}
      sản phẩm là \hl{$g_\theta$}, \no{không phải} nhãn cho tầng b.
    \end{bx}
  \end{column}
  \begin{column}{0.50\textwidth}
    \begin{bx}[title={pha 2 · $N_2 = 200{,}000$ · đường rẻ}]
      \begin{dia}
C0  chẻ + bảng phủ               400
      thiếu mảnh -> THOÁT SỚM
C1  v(rỗng)                      500
C2  g_theta chấm 40 đoạn đơn       0
C3  v(R)                       3.500
C4  v(C)                      10.500
      v_ref = max(v(R), v(C)); tau_q chốt
C5  tìm theo tầng bằng g_theta     0
C6  verification bằng llm THẬT ~7.400
C7  ba nhánh (slide sau)
                              --------
                               ~22.300 token
      \end{dia}
      \vspace{0.2em}
      \dd{\texttt{c2} phải đứng trước \texttt{c5} — hợp lệ vì $g_\theta$ không nhận $\tau_q$}
    \end{bx}
  \end{column}
  \end{columns}
\end{frame}

\begin{frame}[fragile]{$\mathtt{c7}$ — ba nhánh}
  \framesubtitle{thi hành}
  \begin{columns}[t]
  \begin{column}{0.52\textwidth}
    \begin{bx}[title={ba nhánh}]
      \begin{dia}
(1) v(T) >= tau_q, adq phủ hết, tối tiểu
    -> XUẤT T sang tầng b

(2) v(T) < tau_q        T KHÔNG đủ
    -> thử ứng viên kế của beam, tối đa 3
       mỗi lần CHỈ 1 lượt (hết đủ thì dừng,
       không kiểm tối tiểu)
    -> qua: nhánh (1)
       3 lần trượt: nhánh (3)

(3) |T| >= 11 · tỉa kẹt ở >= 6 · thử hết
    -> HÀNG CỨU có phân tầng
      \end{dia}
    \end{bx}
    \begin{keyb}[title={nhánh 2 rẻ hơn tưởng}]
      một ca trượt chỉ tốn \hl{1 lượt} ($1{,}250$ token), một lượt cứu tốn
      $113{,}000$ — \yes{90 lần thử bằng giá một lần cứu}. với $30\%$ qua mỗi lần,
      3 lần vớt $1 - 0{,}7^3 = 66\%$ số ca: $p$ hiệu dụng $0{,}75 \to 0{,}91$.
    \end{keyb}
  \end{column}
  \begin{column}{0.44\textwidth}
    \begin{bx}[title={«đủ mà không tối tiểu» KHÔNG vào hàng cứu}]
      \texttt{c6} đo \emph{trọn} dàn tập con dưới $T$, nên tối tiểu đọc ra ngay,
      0 lượt thêm:
      \begin{dia}
v({a,b,c}) = 0,80 >= tau_q  T đủ
v({a,b})   = 0,78 >= tau_q  con cũng đủ
                   -> T KHÔNG tối tiểu
v({a}) = 0,30  v({b}) = 0,22
                   -> T = {a,b}
      \end{dia}
      \vspace{0.2em}
      ca này xảy ra \emph{thường xuyên}, vì lỗi hay gặp nhất của $g_\theta$ là
      đề xuất tập \no{phình}.
    \end{bx}
    \begin{bx}[title={cứu là gì}]
      chạy \texttt{a1}--\texttt{a4} lên truy vấn đó. $v(\emptyset), v(R), v(C)$ đã có
      trong cache nên còn $113{,}000$ token.
      \vspace{0.15em}

      \dd{$\tau_q$ \emph{tăng} khi cứu — pha 2 vốn đang đo trần thiếu.}
    \end{bx}
  \end{column}
  \end{columns}
\end{frame}

\begin{frame}[fragile]{năm cổng — và chúng nằm trước 75\% chi phí}
  \framesubtitle{thi hành}
  \vspace{0.3em}
  \begin{center}
  \footnotesize
  \begin{tabular}[t]{@{}llrrl@{}}
    \toprule
    & \textbf{giai đoạn} & \textbf{chi phí} & \textbf{tích lũy} & \textbf{cổng chặn sau đó}\\
    \midrule
    \texttt{0}    & hạ tầng · index test có sẵn & \$0 & \$0 &\\
    \texttt{0.5}  & chạy thử bốn cổng trên test index & \$18 & \$18
      & \hl{cổng 0.5} — chặn \$1.100\\
    \texttt{1}    & hiệu chuẩn chuẩn ngoài & \$7 & \$25 & \hl{cổng 1} — $\mathrm{adq}$ đáng tin?\\
    \texttt{2.0}  & index phần pha 1 & \$74 & \$99 &\\
    \texttt{2}    & pha 1 đào nhãn & \$103 & \$202 &\\
    \texttt{3}    & surrogate & \$0 & \$202 & \hl{cổng 2} — vượt gbm?\\
    \texttt{4}    & cổng $p$ trên $N_1$ đủ & \$0 & \$202 & \hl{cổng 3} — chặn \$918\\
    \texttt{5.0}  & index phần pha 2 & \$815 & \$1.017 &\\
    \texttt{5}    & pha 2 $+$ cứu & \$103 & \$1.120 & \hl{cổng 4} $\times\,10$ lô\\
    \bottomrule
  \end{tabular}
  \end{center}

  \vspace{0.6em}
  \begin{keyb}
    tiêu trước \hl{cổng 0.5} là \$18; tiêu trước \hl{cổng 3} là \$202. nếu thiết kế sai
    hoặc $p$ thấp thì \yes{không bao giờ trả \$918 còn lại.}
    \vspace{0.2em}

    \dd{sinh chạy local (qwen3-8b) nên \$0 tiền, nhưng ${\sim}53$ giờ gpu ở 20 sinh/giây —
        ràng buộc thật là \hl{giờ gpu}, không phải tiền. $800+$ của bảng trên là OpenIE
        dựng index cho train, vì index dựng sẵn chỉ có cho \emph{test} set.}
  \end{keyb}
\end{frame}

\begin{frame}[fragile]{những gì chưa chắc — và cổng nào đo nó}
  \framesubtitle{thi hành}
  \begin{center}
  \begin{tabular}[t]{@{}>{\raggedright\arraybackslash}p{0.30\linewidth}
                        >{\raggedright\arraybackslash}p{0.21\linewidth}
                        >{\raggedright\arraybackslash}p{0.41\linewidth}@{}}
    \toprule
    \textbf{con số} & \textbf{trạng thái} & \textbf{ai đo}\\
    \midrule
    độ dài đoạn thật & giả định $250$ tok
      & \texttt{0.b} — sai thì \no{mọi} bảng chi phí sai, và \no{không cổng nào bắt được}\\[0.4em]
    $p$ (tỉ lệ qua verification) & giả định $0{,}75$
      & \texttt{cổng 0.5} đo rẻ, \texttt{cổng 3} đo thật\\[0.4em]
    $45\%$ thoát ở \texttt{a0}/\texttt{c0} & suy từ musique
      & \texttt{2.1} — và chắc chắn \emph{lạc quan}, vì ta cố ý dùng checkpoint chưa fine-tune\\[0.4em]
    $\mathrm{adq}$ sai bao nhiêu & chưa biết
      & \texttt{cổng 1}, đối chiếu \texttt{supporting\_documents}\\[0.4em]
    cỡ rò rỉ nhãn vàng & chưa biết
      & \texttt{1.3} — chênh giữa pool chỉ-bm25 và pool đầy đủ\\
    \bottomrule
  \end{tabular}
  \end{center}

  \vspace{0.5em}
  \begin{keyb}
    hai cổng \yes{miễn phí} phải chạy trước khi cam kết ngân sách:
    \hl{cổng $p$} ($500$ truy vấn pha 1 giữ lại, 0 lượt llm) và
    \hl{chuẩn ngoài} ($500$ truy vấn test, \$7).
    không cổng nào được bỏ qua để «chạy nhanh hơn» — chúng rẻ \emph{chính vì} nằm trước chi phí.
  \end{keyb}
\end{frame}

\begin{frame}{tóm lại}
  \vspace{0.4em}
  \begin{columns}[t]
  \begin{column}{0.48\textwidth}
    \begin{keyb}[title={ba quyết định thiết kế}]
      \textbf{1}\; «đủ» có \emph{hai} thành phần. $v$ đo \emph{trả lời đúng},
      $\mathrm{adq}$ đo \emph{có bằng chứng}. chỉ dùng $v$ thì nhận cả ca mồi tham số.
      \vspace{0.4em}

      \textbf{2}\; tối tiểu, không phải điểm số. mwc cho cả \emph{họ} lời giải thích;
      cái giá là mất completeness.
      \vspace{0.4em}

      \textbf{3}\; chỉ $T$ qua ranh giới. hai tầng đo trên hai không gian khác nhau,
      nên điểm không mang qua được.
    \end{keyb}
  \end{column}
  \begin{column}{0.48\textwidth}
    \begin{bx}[title={hai thứ còn phải chứng minh}]
      \textbf{$p$} — tỉ lệ truy vấn mà đường rẻ qua được verification. cả bảng chi phí
      treo vào nó, và nó \no{chưa được đo}.
      \vspace{0.4em}

      \textbf{$\mathrm{adq}$} — model nli loại đúng hay loại bừa. chưa biết, và nó đang
      được giao quyền bỏ ${\sim}45\%$ truy vấn.
    \end{bx}
    \begin{bx}[title={bước tiếp theo}]
      \texttt{0.a}--\texttt{0.j} trên test index (đã dựng sẵn, \$0), rồi
      \texttt{cổng 0.5} với \$18.
      \vspace{0.2em}

      \dd{github.com/pqvu0307/minimal\_set\_rag-graph\_learning}
    \end{bx}
  \end{column}
  \end{columns}
\end{frame}

\end{document}
